\documentclass[11pt]{article}
\usepackage[utf8]{inputenc}
\usepackage[T1]{fontenc}
\usepackage{geometry}
\usepackage{hyperref}
\usepackage{booktabs}
\usepackage{amsmath,amssymb}
\geometry{margin=2.5cm}
\hypersetup{colorlinks=true,linkcolor=black,urlcolor=blue,citecolor=black}

\title{Modelo multiobjetivo para cobertura de conceptos en redes narrativas}
\author{Documento técnico de trabajo}
\date{29 de julio de 2026}

\begin{document}
\maketitle

\section{Propósito}

Este documento formaliza el problema que resuelve el módulo de optimización de SIAN. El objetivo no es elegir ``las palabras más frecuentes'', sino seleccionar un subconjunto pequeño de conceptos que cubra unidades narrativas relevantes, conserve alta relevancia semántica y mida el daño estructural que produciría retirar esos conceptos de la red.

\section{Red narrativa}

Sea $G=(V,E)$ una red narrativa. Cada nodo $i \in V$ puede representar un monograma, bigrama, trigrama, actor, fuente, marco o evento extraído del corpus. Cada arista $(i,j)\in E$ representa coaparición, relación semántica o relación estructural entre dos nodos.

Los pesos se normalizan en $[0,1]$:

\[
w_i = \frac{c_i}{\max_{k\in V} c_k}, \qquad
w_{ij} = \frac{c_{ij}}{\max_{(p,q)\in E} c_{pq}},
\]

donde $c_i$ es el conteo del nodo y $c_{ij}$ el conteo o intensidad de la relación. Esta normalización evita mezclar escalas incompatibles entre frecuencia textual y peso de aristas.

\section{Universo de cobertura}

Sea $\mathcal{U}$ el universo de unidades narrativas que deben ser cubiertas. Una unidad puede ser:

\begin{itemize}
    \item un documento;
    \item un marco narrativo;
    \item una comunidad semántica;
    \item un evento;
    \item una combinación fuente-año-tema.
\end{itemize}

Para cada unidad $u\in \mathcal{U}$ se define $S_u\subseteq V$, el conjunto de nodos que representan o explican esa unidad. Una solución selecciona nodos mediante variables binarias:

\[
x_i =
\begin{cases}
1, & \text{si el nodo } i \text{ es seleccionado},\\
0, & \text{en otro caso}.
\end{cases}
\]

La factibilidad se define por:

\[
\sum_{i\in S_u} x_i \geq 1 \qquad \forall u\in \mathcal{U}.
\]

Esto es una variante del problema de cobertura de conjuntos. Es NP-completo en su versión de decisión, por lo que se justifica usar heurísticas y metaheurísticas cuando la red crece.

\section{Aristas removidas}

Para medir impacto estructural se define:

\[
z_{ij} =
\begin{cases}
1, & \text{si la arista } (i,j) \text{ se elimina al retirar el conjunto seleccionado},\\
0, & \text{en otro caso}.
\end{cases}
\]

La arista se elimina si al menos uno de sus extremos fue seleccionado y retirado. Por tanto:

\[
z_{ij} \geq x_i,\qquad z_{ij}\geq x_j,\qquad z_{ij}\leq x_i+x_j.
\]

Esto corrige una confusión importante: no se conserva una arista sólo porque ambos nodos estén seleccionados. El objetivo estructural cuenta las aristas que desaparecerían si se retirara el conjunto cubridor.

\section{Objetivos normalizados}

Se usan tres objetivos comparables:

\[
u_1(x)=1-\frac{\sum_{i\in V}x_i}{|V|},
\]

\[
u_2(x)=\frac{\sum_{i\in V} w_i x_i}{\sum_{i\in V} w_i},
\]

\[
u_3(x)=1-\frac{\sum_{(i,j)\in E}w_{ij}z_{ij}}{\sum_{(i,j)\in E}w_{ij}}.
\]

La interpretación es:

\begin{itemize}
    \item $u_1$: maximizar compacidad; usar pocos nodos.
    \item $u_2$: maximizar relevancia; seleccionar nodos con alto peso narrativo.
    \item $u_3$: maximizar preservación estructural; minimizar el peso de aristas removidas.
\end{itemize}

El problema multiobjetivo es:

\[
\max \; F(x) = (u_1(x),u_2(x),u_3(x))
\]

sujeto a las restricciones de cobertura y definición de $z_{ij}$.

\section{Factibilidad}

Las heurísticas deben usar el mismo criterio de factibilidad:

\begin{enumerate}
    \item entre dos soluciones factibles, usar dominancia de Pareto;
    \item entre una factible y una infactible, preferir la factible;
    \item entre dos infactibles, preferir la de menor violación total:
    \[
    \phi(x)=\sum_{u\in \mathcal{U}}\max\left(0,1-\sum_{i\in S_u}x_i\right).
    \]
\end{enumerate}

Este criterio evita comparar soluciones inviables como si fueran resultados científicos.

\section{Métodos comparables}

Los métodos deben ejecutarse bajo el mismo número de evaluaciones de función objetivo y sobre el mismo grafo:

\begin{center}
\begin{tabular}{p{4cm}p{9cm}}
\toprule
Método & Papel\\
\midrule
Barrido glotón ponderado & Construye soluciones con distintos pesos sobre $u_1,u_2,u_3$.\\
MOEA & Cruza y muta subconjuntos, conserva archivo no dominado.\\
MOSA & Usa recocido multiobjetivo con aceptación por dominancia, temperatura y factibilidad.\\
MMC-MO & Parte de soluciones guía obtenidas por relajaciones lineales y modifica arreglos de nodos mediante memoria adaptativa.\\
\bottomrule
\end{tabular}
\end{center}

\section{Relajaciones lineales}

Antes de ejecutar heurísticas se resuelven relajaciones monoobjetivo para estimar puntos guía:

\[
\max u_1(x),\qquad \max u_2(x),\qquad \max u_3(x),
\]

con $x_i,z_{ij}\in[0,1]$. Estas soluciones no sustituyen el problema entero, pero sirven para construir memoria inicial, estimar rangos y definir puntos de referencia para hipervolumen.

\section{Evaluación estadística}

El hipervolumen, la distancia generacional inversa y la dispersión deben calcularse sólo sobre el conjunto no dominado de cada método. Cada método debe ejecutarse al menos $k=10$ veces. Se reportan:

\begin{itemize}
    \item media, mediana, moda, mínimo, máximo y varianza;
    \item intervalo de confianza bootstrap al 95\%;
    \item prueba de Wilcoxon pareada cuando se comparan métodos bajo las mismas semillas;
    \item frente no dominado en tres dimensiones y proyecciones bidimensionales auxiliares.
\end{itemize}

\section{Alcance}

El modelo no demuestra por sí mismo que un concepto sea culturalmente central. Produce candidatos estructuralmente relevantes que deben interpretarse con análisis humano, revisión de fuente y lectura contextual. Esa combinación es coherente con la metodología de análisis narrativo asistido por IA referida en Contactos \cite{contactos625}.

\begin{thebibliography}{9}
\bibitem{contactos625}
Mora-Gutiérrez, R. A. y colaboradores. Artículo en \emph{Contactos}. Disponible en: \url{https://contactos.izt.uam.mx/index.php/contactos/article/view/625}.
\end{thebibliography}

\end{document}

